\documentclass{article}
\usepackage{amsmath,amssymb,amsthm}
\usepackage{algorithm}
\usepackage{algpseudocode}
\usepackage{enumitem}
\usepackage{hyperref}
\usepackage{bm}

\newtheorem{theorem}{Theorem}
\newtheorem{lemma}{Lemma}
\newtheorem{assumption}{Assumption}

\begin{document}

\section*{Gradient Descent--Ascent (GDA)}

\section*{Mathematical Formulation}
Let $L:\mathbb{R}^{d_x}\times\mathbb{R}^{d_y}\rightarrow\mathbb{R}$ be a differentiable function.
We are interested in the saddle‑point problem  

\[
\min_{y\in\mathbb{R}^{d_y}}\;\max_{x\in\mathbb{R}^{d_x}} L(x,y).
\]

Denote $z=(x,y)$ and the partial gradients
\[
\nabla_x L(x,y)\in\mathbb{R}^{d_x},\qquad 
\nabla_y L(x,y)\in\mathbb{R}^{d_y}.
\]

For a stepsize $\eta>0$ the **Gradient Descent–Ascent (GDA)** iteration is  

\[
\boxed{
\begin{aligned}
x_{t+1} &= x_{t} + \eta \,\nabla_x L(x_t,y_t) ,\\[2mm]
y_{t+1} &= y_{t} - \eta \,\nabla_y L(x_t,y_t) .
\end{aligned}}
\tag{1}
\]

Both variables use the same stepsize $\eta$; $x$ moves in the direction of ascent,
$y$ moves in the direction of descent.

\section*{Pseudocode}
\begin{algorithm}[H]
\caption{Gradient Descent–Ascent (GDA)}\label{alg:GDA}
\begin{algorithmic}[1]
\Require Initial point $(x_0,y_0)$, stepsize $\eta>0$, maximum iterations $T$.
\For{$t=0,1,\dots,T-1$}
    \State Compute partial gradients 
    \[
    g^x_t \gets \nabla_x L(x_t,y_t),\qquad 
    g^y_t \gets \nabla_y L(x_t,y_t).
    \]
    \State Update the primal variable (ascent)  
    \[
    x_{t+1} \gets x_t + \eta\, g^x_t .
    \]
    \State Update the dual variable (descent)  
    \[
    y_{t+1} \gets y_t - \eta\, g^y_t .
    \]
\EndFor
\State \Return $\{(x_t,y_t)\}_{t=0}^{T}$.
\end{algorithmic}
\end{algorithm}

\section*{Dry Run Example}
Consider the simple scalar function  

\[
f(w)= (w-3)^2 ,\qquad w\in\mathbb{R}.
\]

Here there is only a primal variable $x\equiv w$ and no dual variable; the
update reduces to gradient descent (a special case of (1) with $y$ absent).

\[
\nabla f(w)=2(w-3).
\]

Initialize $w_0=0$, stepsize $\eta=0.1$, and run three iterations.

\begin{center}
\begin{tabular}{c|c|c|c}
$t$ & $w_t$ & $\displaystyle g_t=\nabla f(w_t)$ & $f(w_t)$\\ \hline
0 & $0$ & $2(0-3)=-6$ & $(0-3)^2=9$\\
1 & $w_1 = 0 + 0.1(-6)= -0.6$ & $2(-0.6-3)=-7.2$ & $(-0.6-3)^2 = 12.96$\\
2 & $w_2 = -0.6 + 0.1(-7.2)= -1.32$ & $2(-1.32-3)=-8.64$ & $(-1.32-3)^2 = 18.6624$\\
3 & $w_3 = -1.32 + 0.1(-8.64)= -2.184$ & $2(-2.184-3)=-10.368$ & $(-2.184-3)^2 = 26.8610$
\end{tabular}
\end{center}

The table shows the gradient evaluation, the update according to (1), and the
objective value after each step.

\newpage

\section*{Assumptions}
\begin{assumption}[Smoothness]\label{as:smooth}
There exists $L>0$ such that for all $z=(x,y)$ and $z'=(x',y')$,
\[
\big\|\nabla L(z)-\nabla L(z')\big\| \le L\|z-z'\|.
\]
Equivalently,
\[
L(z')\le L(z)+\langle\nabla L(z),z'-z\rangle+\frac{L}{2}\|z'-z\|^{2}.
\tag{2}
\]
\end{assumption}

\begin{assumption}[Existence of a Saddle Point]\label{as:saddle}
There exists $(x^{\star},y^{\star})$ such that
\[
L(x^{\star},y)\le L(x^{\star},y^{\star})\le L(x,y^{\star}),\qquad
\forall\,x\in\mathbb{R}^{d_x},\;y\in\mathbb{R}^{d_y}.
\]
Denote the optimal value $L^{\star}=L(x^{\star},y^{\star})$.
\end{assumption}

\begin{assumption}[Bounded Initial Gap]\label{as:gap}
The initial potential is finite,
\[
\Delta_0:=L(x_0,y_0)-L^{\star}<\infty .
\]
\end{assumption}

All assumptions are deterministic; no stochasticity is involved.

\section*{Main Convergence Theorem}
\begin{theorem}[Convergence of GDA]\label{thm:main}
Assume \ref{as:smooth}--\ref{as:gap} hold and the stepsize satisfies  

\[
0<\eta\le \frac{1}{L}.
\]

Define the averaged gradient norm  

\[
\mathcal{G}_T \;:=\; \frac{1}{T}\sum_{t=0}^{T-1}
\Bigl(\|\nabla_x L(x_t,y_t)\|^{2}+\|\nabla_y L(x_t,y_t)\|^{2}\Bigr).
\]

Then for any $T\ge1$,
\[
\boxed{\;
\mathcal{G}_T \le \frac{2L}{\eta T}\,\Delta_0 .
\;}
\tag{3}
\]
Consequently $\mathcal{G}_T = O\!\left(\frac{1}{T}\right)$.
\end{theorem}

\section*{Theoretical Properties}
\begin{itemize}[leftmargin=*]
\item \textbf{Stability.} The condition $\eta\le 1/L$ guarantees that the
potential $L(x_t,y_t)$ never increases beyond the initial value $L(x_0,y_0)$.
\item \textbf{Hyper‑parameter sensitivity.} The bound (3) is inversely
proportional to $\eta$; larger stepsizes (still respecting $\eta\le1/L$) give a
tighter constant but may violate the stability condition.
\item \textbf{Comparison to baselines.}
    \begin{enumerate}
        \item For convex–concave $L$, the same $O(1/T)$ rate holds for the
        extragradient method, but GDA is simpler to implement.
        \item In the non‑convex–concave regime, (3) is the strongest guarantee
        obtainable without additional regularisation (e.g., proximal terms).
    \end{enumerate}
\end{itemize}

\section*{Discussion}
The theorem shows that GDA drives the norm of the stacked gradient
$(\nabla_x L,\nabla_y L)$ to zero at a sublinear rate.  Since a stationary point of
the saddle‑point problem satisfies $\nabla_x L=0$ and $\nabla_y L=0$, the result
implies that the iterates approach a first‑order stationary saddle point in the
average sense.

\newpage

\section*{Preliminaries (Definitions and Lemmas)}
\begin{definition}[Potential function]
\[
\Phi_t \;:=\; L(x_t,y_t)-L^{\star}.
\]
\end{definition}

\begin{lemma}[Descent–Ascent Lemma]\label{lem:da}
Under Assumption~\ref{as:smooth} and the update (1) with $0<\eta\le 1/L$,
\[
\Phi_{t+1}\le \Phi_t
      -\frac{\eta}{2}\bigl(\|\nabla_x L(x_t,y_t)\|^{2}
                         +\|\nabla_y L(x_t,y_t)\|^{2}\bigr).
\tag{4}
\]
\end{lemma}
\begin{proof}
Apply the smoothness inequality (2) twice, once for the $x$‑update (ascent)
and once for the $y$‑update (descent).  For brevity we write $g^x_t:=\nabla_x
L(x_t,y_t)$ and $g^y_t:=\nabla_y L(x_t,y_t)$.  

\textbf{Ascent in $x$:}
\[
\begin{aligned}
L(x_{t+1},y_t)
&\le L(x_t,y_t) + \langle g^x_t, x_{t+1}-x_t\rangle
                 +\frac{L}{2}\|x_{t+1}-x_t\|^{2}\\
&= L(x_t,y_t) + \langle g^x_t, \eta g^x_t\rangle
                 +\frac{L}{2}\eta^{2}\|g^x_t\|^{2}\\
&= L(x_t,y_t) + \eta\|g^x_t\|^{2}
                 +\frac{L}{2}\eta^{2}\|g^x_t\|^{2}.
\end{aligned}
\tag{5}
\]

\textbf{Descent in $y$:}
\[
\begin{aligned}
L(x_{t+1},y_{t+1})
&\le L(x_{t+1},y_t) + \langle g^y_t, y_{t+1}-y_t\rangle
                 +\frac{L}{2}\|y_{t+1}-y_t\|^{2}\\
&= L(x_{t+1},y_t) - \eta\|g^y_t\|^{2}
                 +\frac{L}{2}\eta^{2}\|g^y_t\|^{2}.
\end{aligned}
\tag{6}
\]

Combine (5) and (6):
\[
\begin{aligned}
L(x_{t+1},y_{t+1})
&\le L(x_t,y_t) + \eta\|g^x_t\|^{2}
               -\eta\|g^y_t\|^{2}
               +\frac{L}{2}\eta^{2}\bigl(\|g^x_t\|^{2}
                                        +\|g^y_t\|^{2}\bigr).
\end{aligned}
\tag{7}
\]

Subtract $L^{\star}$ from both sides and rearrange:
\[
\Phi_{t+1}
\le \Phi_t
    -\eta\bigl(\|g^x_t\|^{2}+\|g^y_t\|^{2}\bigr)
    +\frac{L}{2}\eta^{2}\bigl(\|g^x_t\|^{2}+\|g^y_t\|^{2}\bigr).
\tag{8}
\]

Because $\eta\le 1/L$, we have $1-\frac{L}{2}\eta \ge \frac{1}{2}$.
Multiplying (8) by this inequality yields (4).  ∎
\end{proof}

\section*{Main Proof (Step‑by‑Step Derivation)}
We prove Theorem~\ref{thm:main}.  Throughout the proof we use the notation
$g_t^x:=\nabla_x L(x_t,y_t)$, $g_t^y:=\nabla_y L(x_t,y_t)$.

\begin{enumerate}
\item[Step 1.]  Sum Lemma~\ref{lem:da} over $t=0,\dots,T-1$:
\[
\sum_{t=0}^{T-1}\Phi_{t+1}
\le \sum_{t=0}^{T-1}\Phi_t
   -\frac{\eta}{2}\sum_{t=0}^{T-1}
    \bigl(\|g_t^x\|^{2}+\|g_t^y\|^{2}\bigr).
\tag{9}
\]

\item[Step 2.]  The left‑hand side telescopes:
\[
\sum_{t=0}^{T-1}\Phi_{t+1}
= \sum_{t=1}^{T}\Phi_{t}
= \Bigl(\sum_{t=0}^{T-1}\Phi_{t}\Bigr)-\Phi_0+\Phi_T .
\tag{10}
\]

\item[Step 3.]  Substitute (10) into (9) and cancel the common sum:
\[
\Phi_T - \Phi_0
\le -\frac{\eta}{2}\sum_{t=0}^{T-1}
    \bigl(\|g_t^x\|^{2}+\|g_t^y\|^{2}\bigr).
\tag{11}
\]

\item[Step 4.]  Rearrange (11) to isolate the gradient sum:
\[
\sum_{t=0}^{T-1}
    \bigl(\|g_t^x\|^{2}+\|g_t^y\|^{2}\bigr)
\le \frac{2}{\eta}\bigl(\Phi_0-\Phi_T\bigr)
\le \frac{2}{\eta}\,\Delta_0,
\tag{12}
\]
where the last inequality uses $\Phi_T\ge0$ because $L^{\star}$ is the global
minimum in $y$ and maximum in $x$.

\item[Step 5.]  Divide both sides of (12) by $T$ and recall the definition of
$\mathcal{G}_T$:
\[
\mathcal{G}_T
= \frac{1}{T}\sum_{t=0}^{T-1}
    \bigl(\|g_t^x\|^{2}+\|g_t^y\|^{2}\bigr)
\le \frac{2}{\eta T}\,\Delta_0 .
\tag{13}
\]

\item[Step 6.]  Since $L$ is $L$‑smooth, $\Delta_0\le L\|z_0-z^{\star}\|^{2}/2$,
which yields an alternative constant $2L\|z_0-z^{\star}\|^{2}/\eta T$.
Thus (13) is equivalent to the bound (3) stated in the theorem.
\end{enumerate}
∎

\section*{Rate Derivation (With Constants)}
From (13) we have explicitly
\[
\mathcal{G}_T \le \underbrace{\frac{2L}{\eta}}_{\text{constant }C}
               \frac{\|z_0-z^{\star}\|^{2}}{T},
\]
so the rate is $O(1/T)$.  Choosing the maximal admissible stepsize
$\eta = 1/L$ minimizes the constant $C$ and gives
\[
\mathcal{G}_T \le \frac{2L^{2}}{T}\|z_0-z^{\star}\|^{2}.
\]

\section*{Special Cases}
\begin{itemize}
\item \textbf{Convex–concave $L$.}  When $L$ is convex in $x$ and concave in $y$,
the stationary point is a global saddle point; the same $O(1/T)$ bound holds,
and the iterates converge to the unique saddle.
\item \textbf{Pure gradient descent (no $y$).}
If $L$ depends only on $x$, (1) reduces to $x_{t+1}=x_t-\eta\nabla L(x_t)$.
Lemma~\ref{lem:da} becomes the standard descent lemma, and (3) recovers the
classical $O(1/T)$ rate for smooth non‑convex optimization.
\item \textbf{Extragradient method.}
Adding an extra extrapolation step improves the constant but does not change
the $O(1/T)$ asymptotic rate under the same assumptions.
\end{itemize}

\end{document}